Los caudales, en \(\mathrm{L/min}\), satisfacen los intervalos marginales garantizados
\[
Q_1\in\left[100\log\frac{15.9}{8.1},\;100\log\frac{16.1}{7.9}\right]
\subset[67.445,71.196],
\]
\[
Q_2\in\left[100\log\frac{7.9}{2.1},\;100\log\frac{8.1}{1.9}\right]
\subset[132.492,145.002].
\]
Los intervalos exactos son disjuntos: no cabe un único caudal. Tampoco cabe un desplazamiento constante, pues los tres desplazamientos registrados pertenecen, respectivamente, a
\[
[-0.05,0.05],\quad[0.55,0.65],\quad[0.15,0.25]\quad\mathrm{mg/L}.
\]
No queda determinada la evolución del desplazamiento entre pares ni después del último; los intervalos de los dos caudales comparten la concentración en \(t=1\) y no constituyen una caja de valores conjuntamente independientes.
